\documentclass[sigplan,10pt]{acmart}

\setcopyright{none}
\renewcommand\footnotetextcopyrightpermission[1]{}
\settopmatter{printfolios=true}
\pagestyle{plain}

\acmConference[ATC '26]{2026 ACM SIGOPS Annual Technical Conference}{November 15--18, 2026}{Hyatt Hotel, Shatin, Hong Kong}
\acmYear{2026}
\acmISBN{}
\acmDOI{}

\title[ActPlane]{ActPlane: Letting Agents Harness Their Own System-Level Behavior with eBPF}

\author{Paper \#XXX}
\affiliation{%
  \institution{Anonymous Submission}
  \city{}
  \country{}}

\begin{abstract}
AI coding agents require harnesses to apply behavioral policies
for effective and safe operation: e.g. do not leak secrets, run tests before committing.
Our empirical study of 64 projects shows that while project harness
instructions are declared in natural language, 64\%
consist of behavioral policies.

Yet existing agent harnesses fail to connect intent-level declarations
to deterministic system-level policies, leaving a \emph{semantic gap} and
making it difficult to ensure compliance or follow rules.
While 83\% of directives involve system-level behavior,
prompt constraints operate at the probabilistic intent level.
81\% of projects require cross-event state tracking, but tool-call
guardrails lose track of system-level actions.
74\% of system-level rules require intent-level project or
task context to evaluate, while current OS-level
mechanisms expect pre-defined static
policy and return only system events without semantic context.

To address this, we argue the policy engine should be exposed as a
programmable interface, allowing agents
to dynamically define and adapt their policy
for their own system-level actions
based on their intent context.
We realize this in ActPlane, which allows agents to declare cross-event
behavior as labeled information-flow control (IFC) policies under 
a layered policy model,
observing and enforcing the agents' actions across process, file, and
network boundaries with semantic feedback.
We implement ActPlane with eBPF and evaluate on 607 policies from the
empirical study, showing that ActPlane improves policy compliance by
XX\% over existing enforcement mechanisms and end-to-end task
completion by XX~pp, with $<$X~\textmu s per-syscall overhead.
\end{abstract}

\begin{CCSXML}
<ccs2012>
 <concept>
  <concept_id>10010520.10010553.10010562</concept_id>
  <concept_desc>Computer systems organization~Embedded and cyber-physical systems</concept_desc>
  <concept_significance>300</concept_significance>
 </concept>
 <concept>
  <concept_id>10002978.10003022.10003023</concept_id>
  <concept_desc>Security and privacy~Software and application security</concept_desc>
  <concept_significance>300</concept_significance>
 </concept>
 <concept>
  <concept_id>10011007.10010940.10010941.10010949</concept_id>
  <concept_desc>Software and its engineering~Software safety</concept_desc>
  <concept_significance>300</concept_significance>
 </concept>
</ccs2012>
\end{CCSXML}

\ccsdesc[300]{Computer systems organization~Embedded and cyber-physical systems}
\ccsdesc[300]{Security and privacy~Software and application security}
\ccsdesc[300]{Software and its engineering~Software safety}

\keywords{AI agents, eBPF, BPF-LSM, labeled information-flow control, information-flow control, provenance, semantic feedback}

\begin{document}
\maketitle

\section{Introduction}

% ¶1: Context — agents are powerful and widely deployed
AI coding agents have become powerful, widely deployed tools for software
development. They operate as long-running processes with access to shells,
source trees, package managers, and external APIs, autonomously writing code,
running tests, and managing infrastructure across extended
sessions~\cite{agentsight}.

% ¶2: Agents need harnesses to enforce behavioral contracts (consensus, state quickly)
As agents gain autonomy, projects govern them with behavioral contracts---do
not push directly to \texttt{main}; run tests before committing; never expose
secrets---encoded in instruction files such as \texttt{CLAUDE.md} and
\texttt{AGENTS.md}~\cite{chatlatanagulchai2025manifests,lulla2026agentsmd}.
Because LLM compliance is inherently probabilistic, deterministic enforcement
cannot live inside the model; it requires an external
harness~\cite{langchain-harness,fowler-harness}. The question is not whether
agents need harnesses, but where the harness should enforce and how it should
feed back.

% ¶3: Gap — single-level enforcement leaves an intent–behavior gap
Existing agent harness enforcement operates at a single level, leaving an
\emph{intent--behavior gap}. An agent's execution spans three abstraction
levels: \emph{intent} (what the agent declares it will and will not do),
\emph{action} (the tool calls it issues through its runtime), and
\emph{behavior} (the actual OS operations---\texttt{open}, \texttt{execve},
\texttt{connect}---that
result)~\cite{agentsight}. The mapping from action to behavior is many-to-many
and opaque: a single \texttt{run\_command("make")} can trigger hundreds of
system calls, and the same \texttt{git~push} can be reached through a direct
tool call, a \texttt{bash~-c} string, a Python subprocess, or a compiled
helper. Intent-level enforcement (prompt instructions) is probabilistic---the
agent may forget. Action-level enforcement (tool-call guards such as AgentSpec
and Progent~\cite{wang2025agentspec,shi2025progent}) is deterministic but
bypassable: it checks tool calls, not the OS behavior they produce.
Behavior-level enforcement (syscall filters, sandboxes) is unbypassable but
returns opaque \texttt{-EPERM} failures disconnected from intent. No single
level bridges the gap.

% ¶4: Insight — contracts are cross-object information flow
Real behavioral contracts frequently require cross-object information-flow
tracking that no single enforcement level provides. A contract like ``never
expose secrets to the network'' is not a predicate on one system call---it
requires knowing that the current process previously read a secret file and
now attempts a network connection. A contract like ``run tests before
committing'' requires knowing that a test process executed before a commit
process in the same session. A corpus study of instruction files from
50~popular projects confirms that such cross-object, cross-temporal contracts
are pervasive, yet today's agent
harnesses~\cite{langchain-harness,fowler-harness} keep their enforcement and
feedback components entirely at the application layer, where they cannot track
information flow across process, file, and network boundaries.

% ¶5: System — ActPlane bridges intent and behavior
ActPlane bridges intent and behavior by compiling agent-declared contracts into
OS-level labeled information-flow enforcement with corrective feedback. Agents
and developers declare contracts in a compact DSL---voluntary confinement,
analogous to \texttt{pledge()} in OpenBSD, where a program restricts its own
capabilities. ActPlane compiles these declarations into labeled
information-flow rules and loads them into an eBPF/BPF-LSM backend. Named
labels propagate across process, file, and network objects at kernel hooks,
encoding execution history as checkable state: when a process reads a file
labeled \texttt{SECRET}, the process acquires that label; when it forks, the
child inherits it. When a labeled subject violates a rule, ActPlane returns
corrective feedback to the agent---which contract was violated, why, and what
alternative path satisfies it---not \texttt{Permission~denied} but a reminder
of the agent's own declared intent.

% ¶6: Positioning — not new mechanism, unification at OS level
ActPlane is not a new enforcement mechanism but a unification at the OS level
of capabilities previously split across layers. CamQuery showed that
cross-channel label propagation and enforcement can be done in the
kernel~\cite{pasquier2018camquery}; AgentSpec demonstrated corrective feedback
for agents at the tool-call layer~\cite{wang2025agentspec}; FIDES and CaMeL
let agents participate in their own governance through application-layer
information-flow control~\cite{costa2025fides,debenedetti2025camel}. ActPlane
brings these together on the modern eBPF/BPF-LSM substrate, where contracts
survive context loss and cannot be bypassed below the tool API. It is not a
replacement for the application-layer harness but a complement---the
enforcement component that must reach below the tool API to the OS.

% ¶7: Results (TBD — fill with headline numbers when available)
% TODO: Add headline results here once evaluation is complete.
% Example: "ActPlane catches all five bypass paths that tool-layer guards miss,
% improves agent task-completion rate from X% to Y% under corrective feedback,
% incurs Z µs per-event overhead, and produces no false positives on N benign
% workloads."

% ¶8: Contributions
This paper makes three contributions.

\begin{enumerate}
\item It defines the \emph{intent--behavior gap} for AI agents and grounds it
with a corpus study of instruction files from 50~popular projects, showing
that cross-object information-flow contracts are pervasive in real agent
instruction files.

\item It presents ActPlane, a system that compiles agent-declared behavioral
contracts into eBPF/BPF-LSM labeled information-flow enforcement with
corrective feedback, propagating labels across process, file, and network
objects and reconnecting behavior-level violations to intent-level
remediation.

\item It evaluates ActPlane on bypass coverage across five tool paths, agent
recovery rate under four feedback conditions, false positives under benign
workloads, and per-event overhead, comparing against action-level and
behavior-level baselines.
\end{enumerate}

\section{Background and Motivation}
\label{sec:motivation}

\subsection{Agent Harnesses and Behavioral Contracts}

An agent harness is the non-model infrastructure that mediates between an LLM
agent and its execution
environment~\cite{langchain-harness,fowler-harness}. Harness components
include tool dispatch, memory management, orchestration, enforcement, and
feedback. This paper focuses on the enforcement and feedback components:
the mechanisms that ensure an agent's behavioral contracts hold at runtime.

Behavioral contracts are constraints that govern an agent's observable
effects rather than its internal reasoning. Projects encode them in
instruction files such as \texttt{CLAUDE.md} and
\texttt{AGENTS.md}~\cite{chatlatanagulchai2025manifests,lulla2026agentsmd}:
do not push directly to \texttt{main}; run tests before committing; never
expose secrets to the network. These contracts differ from coding-style
directives (e.g., ``prefer \texttt{const} over \texttt{let}'') in that they
produce observable effects at the system-call boundary.

\subsection{The Intent--Behavior Gap}

An agent's execution spans three abstraction levels.

\textbf{Intent} is the agent's own goals and self-constraints, i.e., what it
declares it will and will not do. 
% In ActPlane's design, the agent
% formalizes its intent in a compact DSL, just as a program formalizes its
% capability intent via \texttt{pledge()}~\cite{openbsd-pledge}. ActPlane
% does not infer intent from the agent's behavior (cf.\ AgentSight's
% passive observation~\cite{agentsight}); it receives intent as an active
% declaration that the agent itself writes and maintains.
\textbf{Action} is the tool calls the agent issues through its runtime
(\texttt{read\_file}, \texttt{run\_command}, \texttt{edit\_file}), mediated
by the agent framework.

\textbf{Behavior} is the actual OS-level operations that result:
\texttt{open}, \texttt{execve}, \texttt{connect}, and \texttt{fork}.

The mapping from action to behavior is many-to-many and opaque. A single
\texttt{run\_command("make")} triggers hundreds of system calls. The same
\texttt{git push} can be reached through a direct tool call, a
\texttt{bash~-c} string, a Python \texttt{subprocess}, or a compiled helper
binary. This opacity creates an \emph{intent--behavior gap}: contracts
declared at the intent level cannot be reliably enforced at a single
intermediate level.

Enforcement placed at the intent level (prompt instructions) is
probabilistic: LLM compliance degrades with context length and task
complexity~\cite{liu2024lost}. Enforcement placed at the action level (tool-call
guards~\cite{wang2025agentspec,shi2025progent}) is deterministic but
bypassable: it intercepts tool calls, not the OS behavior they produce.
Enforcement placed at the behavior level (syscall filters,
sandboxes~\cite{pasquier2018camquery,ciliumtetragon}) is complete but
returns opaque \texttt{-EPERM} failures with no connection to the agent's
declared contracts. Each level addresses a different part of the stack;
none spans the full gap from intent to behavior.

A cooperative agent does not intentionally bypass tool-layer guards, but
the opacity of the action-to-behavior mapping causes unintentional bypass
in routine operation. Three representative scenarios illustrate this.
(1)~An agent calls \texttt{run\_command("npm~test")}; the test harness
internally invokes \texttt{git~push} to a staging server. The tool-call
guard sees \texttt{npm~test} (a permitted action), while the kernel
observes the forbidden \texttt{git~push} behavior.
(2)~An agent calls \texttt{read\_file(".env")} to inspect a configuration
value, then later calls \texttt{run\_command("curl~api.example.com")}.
Each tool call is individually permitted, but the combination constitutes
secret exfiltration, a cross-object violation visible only to an enforcer
that tracks information flow from file read to network connect.
(3)~An agent writes a Python script containing
\texttt{subprocess.run(["git",~"push",~"--force"])} and executes it. The
tool-call guard sees \texttt{run\_command("python~script.py")}; the kernel
observes the nested \texttt{execve("git",~["push",~"--force"])}.

\subsection{Threat Model}

ActPlane targets a \emph{cooperative but probabilistic} agent. The agent is
not adversarial: it does not attempt to defeat the enforcement mechanism
through obfuscation or exploitation. However, because LLM inference is
inherently stochastic, the agent may violate any given contract with
non-zero probability at each step, regardless of context length or prompt
design. Over a long session with many steps, the cumulative probability of
at least one violation approaches certainty.

The trust boundary is the kernel. The eBPF programs and BPF maps that
propagate labels and check rules run in kernel space and cannot be modified
by the agent. The policy DSL and feedback formatting run in a trusted
userspace process (\texttt{actplane run}) that is not part of the agent's
process tree. The monitored domain is the agent's entire process tree,
including all descendants spawned through tools, shells, or subprocesses.

ActPlane does not defend against an agent that deliberately exploits kernel
vulnerabilities or manipulates the eBPF loader. Such attacks require
\texttt{CAP\_BPF} or root privileges that the agent does not hold.
Side-channel inference (deducing policy rules from \texttt{-EPERM}
patterns), deliberate label pollution (triggering over-tainting to cause
denial of service), and feedback-oracle attacks (using remediation strings
to restructure an approach toward a forbidden goal through compliant steps)
all require adversarial intent and are outside this threat model.


\section{Corpus Study}
\label{sec:corpus}

This section characterizes behavioral contracts in real agent instruction
files and identifies the subset that requires cross-object information-flow
enforcement.

\subsection{Methodology}

We collected \texttt{CLAUDE.md} and \texttt{AGENTS.md} files from 50
popular open-source projects, selected by GitHub star count and filtered
to exclude documentation-only repositories. From these files we extracted
866 candidate imperative statements using keyword matching (``never,''
``must,'' ``do not,'' ``before,'' ``only'') and manually coded each
statement along seven dimensions: (D1)~coding style vs.\ behavioral
contract, (D2)~contract category, (D3)~observable at the syscall boundary,
(D4)~expressible in the ActPlane DSL, (D5)~enforceable without false
positives under normal use, (D6)~robust to naming and path variation, and
(D7)~temptable (likely to be violated by an agent in a realistic task).
Two annotators independently coded the full set; we report Cohen's
$\kappa$ for inter-rater reliability.

\subsection{Per-Action vs.\ Cross-Object Contracts}

Behavioral contracts fall into two categories that differ in enforcement
requirements.

\textbf{Per-action contracts} constrain a single operation independently of
execution history. Examples include ``do not execute \texttt{rm~-rf}'' and
``do not run \texttt{git branch}.'' These contracts can be enforced by
matching a single system call or tool call against a static rule.

\textbf{Cross-object contracts} constrain an operation based on prior
interactions with other objects. Examples include ``a process that has read
\texttt{.env} must not connect to an external endpoint'' (data-flow
contract) and ``commit only after tests have passed'' (temporal contract).
These contracts require tracking information flow or execution ordering
across process, file, and network boundaries.

% TODO: Insert table with category breakdown and percentages once coding
% is complete.

\subsection{Implications for Enforcement}

Per-action contracts can be enforced by per-event matching at any level:
tool-call guards, seccomp filters, or eBPF kprobes. Cross-object contracts
cannot, because they depend on state accumulated across multiple operations
and objects. Enforcing them requires a mechanism that (1)~tracks provenance
across process, file, and network boundaries, (2)~encodes that provenance
as checkable state at each enforcement point, and (3)~connects violations
back to the declared contract for corrective feedback.

Labeled information-flow control satisfies all three requirements. Labels
encode cross-object provenance as per-node bitmasks, propagation rules
maintain the invariant across fork, exec, read, write, and connect
operations, and source-to-sink rules check the accumulated state at each
enforcement point. The ActPlane DSL compiles cross-object contracts
directly into this model.

\section{Design}
\label{sec:design}

ActPlane is a programmable OS-level control plane for agent harnesses. It
compiles behavioral policies into labeled information-flow rules at OS
kernel hooks, observing and enforcing agent behavior and returning semantic
feedback when a rule is matched. Figure~\ref{fig:architecture} summarizes
the pipeline.

\begin{figure}[t]
\begin{verbatim}
actplane.yaml
  policy: |
        |
        v
Rust compiler ------> struct taint_config
        |                    |
        v                    v
 actplane run -----> eBPF loader + BPF maps
        |                    |
        v                    v
  target command <--- kernel hooks
        |
        v
.actplane/last-violation.txt
\end{verbatim}
\caption{ActPlane pipeline. The Rust compiler lowers the policy DSL into a
fixed-size binary configuration. The eBPF loader writes it into read-only
data, attaches to kernel hooks, and reports rule matches as NDJSON. The runner
maps rule identifiers back to human-readable feedback.}
\Description{ASCII diagram of policy compilation, eBPF loading, target
command execution, and feedback output.}
\label{fig:architecture}
\end{figure}

\subsection{Policy Model}

ActPlane policies are object-level information-flow rules over three node
types: processes (identified by pid and command name), files (keyed by
path hash), and endpoints (keyed by IPv4 address and mask). The model
operates at object granularity rather than byte
granularity~\cite{kemerlis2012libdft}, trading precision for low-overhead
whole-system coverage.

\subsubsection{Labels and Propagation}

Each node carries a 64-bit label bitmask. A \emph{source} declaration
introduces labels: an exec source labels processes that execute a matching
binary, a file source labels files matching a path pattern, and an endpoint
source labels network endpoints matching an address range. Labels propagate
through fixed transfer functions:

Let $L(x)$ denote the label bitmask of node $x$. Labels propagate through
fixed transfer functions:

\begin{itemize}
\item \textbf{fork}$(p \to c)$: $L(c) := L(c) \cup L(p)$.
\item \textbf{exec}$(p, f)$: $L(p) := L(p) \cup L(f) \cup S(f)$, where
$S(f)$ is the set of source labels matching $f$.
\item \textbf{read}$(p, f)$: $L(p) := L(p) \cup L(f)$.
\item \textbf{write}$(p, f)$: $L(f) := L(f) \cup L(p)$.
\item \textbf{connect}$(p, e)$: $L(e) := L(e) \cup L(p)$.
\end{itemize}

All transfer functions are monotonic: labels accumulate and are never
removed by propagation. The only mechanism that removes a label is an
explicit \emph{declassify} directive in the DSL, which clears specified
labels when the agent executes a designated gate binary (e.g., a redaction
tool). This monotonicity bounds over-tainting: a node's label set grows
at most to the union of all source labels reachable through observed flow
edges.

These transfer functions maintain a provenance invariant: if a process $p$
carries label $\ell$, then information from some source tagged with $\ell$
has reached $p$ through observed OS-level flow edges.

\subsubsection{Rules and Effects}

A rule contains one or more clauses. Each clause starts with an action verb
(\texttt{notify}, \texttt{block}, or \texttt{kill}) followed by an
operation (\texttt{exec}, \texttt{read}, \texttt{write}, \texttt{unlink},
\texttt{connect}, or \texttt{open}), a target pattern, optional positional
arguments, a Boolean expression over required and forbidden labels, and an
optional condition (target allow-list, lineage gate, or temporal after-gate).

\begin{figure}[t]
\begin{verbatim}
source AGENT = exec "**/claude"
source SECRET = file "**/.env"

rule no-secret-egress:
  block connect endpoint "*" if SECRET
  because "secret-derived data must stay local; run the redactor first"

rule no-agent-git:
  kill exec "git" if AGENT
  because "agent must not invoke git directly"
\end{verbatim}
\caption{Two ActPlane policy rules. The first uses BPF-LSM pre-operation
denial; the second uses harness-level termination. Exec target patterns
without \texttt{/} are implicitly basename-matched.}
\Description{Code listing of two ActPlane rules.}
\label{fig:policy-example}
\end{figure}

Each clause declares its own effect that determines how the kernel
responds to a match.

\textbf{Notify}: the kernel emits a match event and allows the operation
to proceed.

\textbf{Block}: a BPF-LSM hook returns \texttt{-EPERM} before the operation
commits. Block is available only when the \texttt{bpf} LSM is active in the
running kernel; it is not silently downgraded.

\textbf{Kill}: the kernel sends \texttt{SIGKILL} to the matching task. From
a BPF-LSM hook, kill also returns \texttt{-EPERM}; from a tracepoint, kill
is post-operation termination.

When multiple rules match the same event, the strongest effect applies in
the order kill $>$ block $>$ notify.

\subsubsection{Policy Language Scope}

The DSL is not a general-purpose mandatory access control policy. Its
primitives map to recurring agent workflow patterns: ``after tests,''
``through a review tool,'' ``data derived from untrusted input,'' and
``read-only sub-agent.'' The grammar is small so that policy authors can
translate natural-language instructions into source/sink/gate shapes without
security-policy expertise.

\subsection{System Architecture}

\subsubsection{Compiler and Kernel ABI}

The Rust compiler parses the DSL and lowers it into a C-compatible
\texttt{struct taint\_config}. Lowering allocates label bits, expands Boolean
expressions into required and forbidden bitmasks, lowers glob patterns into
exact, prefix, suffix, or wildcard match kinds, and compiles IPv4 patterns
into network/mask pairs. Human-readable rule names and reason strings
(including remediation guidance) remain in Rust metadata; the kernel receives
only fixed-size rule tables and emits integer rule identifiers.

This split keeps the eBPF program verifier-friendly. The kernel program
iterates over fixed-size arrays in read-only data using bounded loops, copies
each rule entry into a stack-local variable before matching, and uses explicit
index bounds checks for argument scanning and pattern comparison.

\subsubsection{Kernel State}

The eBPF backend maintains five BPF hash maps:

\begin{itemize}
\item \texttt{ts\_proc}: pid $\to$ label bitmask and lineage-gate state.
\item \texttt{ts\_root}, \texttt{ts\_sess}: root-process and session-level
temporal (after) gate state.
\item \texttt{ts\_file}: path hash $\to$ label bitmask.
\item \texttt{ts\_endp}: IPv4 address $\to$ label bitmask.
\end{itemize}

BPF map entries are keyed by tid (one entry per Linux thread) for
\texttt{ts\_proc} and by path hash or IPv4 address for object maps.
Updates use atomic BPF helper operations; concurrent reads from different
CPUs observe a consistent label bitmask per entry.

Tracepoint programs handle process lifecycle events (fork, exec, exit) and
post-operation file and connect events. BPF-LSM programs attach to
\texttt{bprm\_check\_security}, \texttt{file\_open}, and
\texttt{socket\_connect} for pre-operation enforcement. When BPF-LSM is not
active, the loader uses tracepoint mode and disables block-effect rules.
In tracepoint mode, hooks fire after the operation commits; a
\texttt{connect} tracepoint observes an already-established connection.
This is why \texttt{block} requires BPF-LSM: only LSM hooks
provide a pre-operation verdict. Tracepoint mode supports \texttt{notify}
(post-operation observation) and \texttt{kill} (post-operation task
termination).

The current implementation supports up to 32 sources, 32 rules, 16
transforms, and 16 gates per policy. Rule evaluation uses
\texttt{bpf\_loop} for bounded iteration, which the eBPF verifier accepts
for policies exceeding 100 rules.

\subsubsection{Unified Event Handling}

Each hook normalizes kernel arguments into a common event structure:
subject pid, object kind, access kind, backend mode, target string,
optional argument vector, and optional IPv4 address. A shared handler
resolves the subject's label bitmask, applies source labels for
read/connect events, evaluates all supported rules, emits at most one
\texttt{TAINT\_VIOLATION} event for the highest-priority matching rule, and
then applies propagation updates.

\subsection{Corrective Feedback}

The kernel match event contains the rule identifier, effect, target
string, pid, and flags indicating whether the operation was blocked or the
task was killed. The Rust runner maps the rule identifier to the
corresponding rule name and reason string, then appends a
structured payload to \texttt{.actplane/last-violation.txt}.

Agent integration uses a post-tool hook. The \texttt{actplane feedback-hook}
subcommand reads new bytes from the feedback file and returns them as
\texttt{additionalContext} to compatible agent
runtimes~\cite{agentsight}. The hook is a read-only relay: it does not
re-evaluate policy in userspace. The kernel is the sole authority for rule
matches; userspace formats the result so the agent can select an alternative
execution path.


\section{Implementation}
\label{sec:implementation}

ActPlane is implemented as a Rust driver and an eBPF program pair. The
system builds a single \texttt{actplane} binary that embeds the compiled
eBPF loader, extracts it at runtime, compiles the project's policy DSL,
and runs a target command under the observation and enforcement harness.

\subsection{Runner and Policy Discovery}

The command \texttt{actplane run -- <cmd>} discovers
\texttt{actplane.yaml} by searching upward from the working directory.
The policy YAML contains a \texttt{policy:~|} DSL block. Before the
target process begins executing, the loader seeds the stopped target's pid
with the \texttt{AGENT} label so that all descendants inherit the agent
identity. The runner prepares the feedback file
(\texttt{.actplane/last-violation.txt}) and a hook state file, exports
their paths to the child environment, waits for the loader to report
readiness, then resumes the target. This sequencing ensures that the
first target operation is observed.

\subsection{BPF Backends}

The eBPF program attaches to process lifecycle tracepoints (fork, exec,
exit), syscall tracepoints for file and connect events, and BPF-LSM hooks
for pre-operation enforcement. At load time, the loader checks whether the
running kernel has the \texttt{bpf} LSM active in
\texttt{/sys/kernel/security/lsm}. If not, it falls back to tracepoint
mode, in which the supported effects are \texttt{notify} and \texttt{kill};
\texttt{block} rules are disabled because there is no pre-operation verdict
path.

\subsection{Limitations}

\textbf{Path identity.} File labels are keyed by FNV-1a path hash rather
than inode/device pairs. Hard links and mount binds can cause the same
file to carry different labels under different paths.

\textbf{Endpoint identity.} Connect matching uses numeric IPv4 addresses.
Hostname, DNS, SNI, and HTTP-layer matching require userspace resolution or
additional instrumentation.

\textbf{Argument predicates.} Positional argument predicates inspect
arguments after exec via the tracepoint argument buffer. Pre-operation
blocking of argument predicates requires an argument cache before
\texttt{bprm\_check\_security}, which is not yet implemented.

\textbf{Label precision.} Labels propagate at object granularity. A process
that reads one byte from a labeled file acquires the full label set. This
can produce over-tainting relative to byte-level
systems~\cite{kemerlis2012libdft}. Section~\ref{sec:evaluation} measures the
false-positive rate and validates that declassification and gate paths
prevent unnecessary harness intervention.

\section{Evaluation}
\label{sec:evaluation}

The evaluation measures five properties: bypass coverage, corrective-feedback
effectiveness, false-positive rate, runtime overhead, and coverage over
real-world contracts.

\subsection{Experimental Setup}

% TODO: Fill in agent versions, kernel version, hardware, policy set.

\subsection{Bypass Coverage (RQ1)}

Does OS-level enforcement detect the same contract violation regardless of
the tool path through which the agent triggers it?

We execute the same forbidden operation through five paths: (1)~direct tool
call, (2)~nested shell string (\texttt{bash~-c}), (3)~Python
\texttt{subprocess}, (4)~compiled helper binary, and (5)~direct
\texttt{execve} from generated C code. For each path, we record whether
ActPlane detects the violation and whether the corrective feedback payload
is delivered to the agent.

% TODO: Insert coverage matrix table.

\subsection{Feedback and Recovery (RQ2)}

Does corrective feedback improve agent task completion and reduce repeated
violations compared with enforcement without feedback?

We run agent tasks under four conditions: (C1)~prompt-only, in which the
contract is stated in the prompt but not enforced; (C2)~audit, in which ActPlane
reports violations but does not fail operations; (C3)~enforcement without
feedback, in which the operation fails with \texttt{-EPERM} or \texttt{SIGKILL}
but no structured reason; (C4)~enforcement with feedback, in which the operation
fails and ActPlane injects the rule name, reason, and remediation into the
agent's context.

The primary comparison is C3 vs.\ C4: the marginal benefit of corrective
feedback. The C1-to-C2 comparison measures how often prompt-only contracts
are violated. The C2-to-C3 comparison isolates the effect of enforcement
from observation.

% TODO: Insert results table.

\subsection{False Positives (RQ3)}

Do benign tasks and sanctioned contract paths (declassification, gate
mediation, temporal gates) proceed without spurious enforcement?

% TODO: Insert false-positive results.

\subsection{Overhead (RQ4)}

What is the per-event and end-to-end cost of label propagation and rule
checking?

We measure per-event latency for fork, exec, open, write, and connect
operations with and without ActPlane attached, and report p50 and p99
latencies. We also measure wall-clock time for representative agent tasks
and BPF map memory consumption as a function of rule count and active
process count.

% TODO: Insert latency table and overhead figures.

\subsection{Corpus Coverage (RQ5)}

How many real-world contracts from the corpus study
(\S\ref{sec:corpus}) are enforceable by ActPlane?

We apply the D1--D7 coding dimensions to each extracted contract and
report the fraction that is observable at the syscall boundary (D3),
expressible in the ActPlane DSL (D4), enforceable without false positives
(D5), and temptable in a realistic agent task (D7). We report inter-rater
agreement (Cohen's $\kappa$) for each dimension.

% TODO: Insert corpus coverage funnel table.


\section{Related Work}
\label{sec:related}

\subsection{In-Kernel Provenance and Information Flow}

PASS introduced provenance as a storage-layer
primitive~\cite{muniswamy2006pass}. Hi-Fi and LPM used LSM hooks to capture
trustworthy whole-system provenance~\cite{pohly2012hifi,bates2015lpm}.
CamFlow extended this to a streaming provenance framework over processes,
files, and sockets~\cite{pasquier2017camflow}.

CamQuery is the closest prior
system~\cite{pasquier2018camquery}. It executes provenance queries inline
in the kernel, propagates labels across process, file, and network edges,
and can deny operations that violate a policy. ActPlane shares this
label-propagation model but differs in three respects: it targets
cooperative AI agents rather than adversarial intrusions, compiles an
agent-oriented source/sink DSL rather than general provenance queries, and
returns corrective feedback to the violating actor. The implementation
differs as well: ActPlane uses the eBPF/BPF-LSM substrate rather than a
loadable kernel module.

TaintDroid, Dytan, libdft, and Panorama established dynamic taint tracking
at runtime, binary-instrumentation, and emulation
layers~\cite{enck2010taintdroid,clause2007dytan,kemerlis2012libdft,yin2007panorama}.
ActPlane operates at object granularity rather than byte granularity,
trading precision for low-overhead whole-system coverage.

\subsection{eBPF and LSM Enforcement}

BPF-LSM, originally proposed as KRSI, allows eBPF programs to attach to
Linux security hooks and return allow/deny
verdicts~\cite{singh2019krsi,linuxbpflsm}. ActPlane uses BPF-LSM as its
pre-operation enforcement backend. Contemporary eBPF enforcement tools
such as Tetragon, Tracee, and eBPF-PATROL provide per-event predicate
matching and, in Tetragon's case, a boolean lineage flag
(\texttt{followChildren}) that propagates across fork/exec. These systems
evaluate per-event predicates or single-channel lineage flags rather than
cross-channel labeled information
flow~\cite{ciliumtetragon,aquasecuritytracee,ghimire2025ebpfpatrol}.

\subsection{Agent Guardrails and Information-Flow Control}

AgentSpec and Progent provide deterministic runtime enforcement over
tool-call names, arguments, and framework
actions~\cite{wang2025agentspec,shi2025progent}. These systems enforce at
the tool-call boundary; ActPlane complements them by enforcing below the
tool API, where effects that bypass the framework are still visible.

FIDES and CaMeL apply typed information-flow control and capability
separation within the agent loop or
scaffolding~\cite{costa2025fides,debenedetti2025camel}. ActPlane applies
the same style of information-flow reasoning at the OS boundary, where
enforcement persists across context windows and cannot be circumvented by
subprocess indirection.

\subsection{Agent Instruction Files}

Recent empirical studies characterize \texttt{CLAUDE.md},
\texttt{AGENTS.md}, and related agent instruction files along dimensions
of content taxonomy, maintenance practices, and effect on task
efficiency~\cite{chatlatanagulchai2025manifests,chatlatanagulchai2025agentreadmes,
santos2025claudeconfig,lulla2026agentsmd}. ActPlane's corpus study
(\S\ref{sec:corpus}) asks a different question: which instructions
constitute cross-object information-flow contracts, and which of those can
be enforced at the OS level?


\section{Conclusion}
\label{sec:conclusion}

ActPlane provides a programmable OS-level control plane for agent harnesses.
It compiles agent-declared behavioral contracts into eBPF/BPF-LSM labeled
information-flow rules, propagates labels across process, file, and endpoint
objects at kernel hooks, and returns kernel-detected violations as semantic
feedback to the agent. By supporting audit, pre-operation denial, and
harness-level termination with structured remediation, ActPlane enables
agents to recover from contract violations rather than fail opaquely.
ActPlane unifies labeled information-flow observation and enforcement with
an agent-oriented contract language and a feedback loop on the eBPF/BPF-LSM
substrate, bridging the intent--behavior gap that single-level harnesses
leave open.


\bibliographystyle{ACM-Reference-Format}
\bibliography{references}

\end{document}
